%!TEX root = these.tex

%\chapterstar{CONCLUSION GÉNÉRALE ET PERSPECTIVES}


\chapter*{Conclusion générale et perspectives}
\mtcaddchapter

\addstarredpart{CONCLUSION GÉNÉRALE ET PERSPECTIVES} % Pour ajouter une partie fictive dans la table des matière
\mtcaddpart

\markboth{Conclusion générale et perspectives}{Conclusion générale et perspectives}



La com



\section*{Contributions de la thèse}

Nous avons mis en évidence, dans la première partie de ce manuscrit, les \textbf{usages actuels en biologie structurale}, coeur d'application de nos travaux de recherche. Ce domaine spécifique de la biologie moléculaire, cherchant à décrypter les structures des protéines afin d'en déduire leurs fonctions, a su évoluer avec les outils informatiques, et ce malgré la part importante des méthodes expérimentales dans la résolution de structures 3d. 
L'informatique a tout d'abord fait partie intégrante du traitement des signaux générés par la majorité des méthodes expérimentales pour être ensuite considérée comme support de méthodes fiables de modélisation de structures protéiques. La modélisation théorique de structures moléculaire reçu d'ailleurs ses plus belles lettres de noblesse en 2013, ses pionniers recevant le prix Nobel de chimie. L'informatique graphique a ensuite constitué une méthode de communication et de diffusion des modèles de structures primordiale pour le monde scientifique, permettant de partager les avancées du domaine au plus grand nombre.

Parmi les nombreuses approches théoriques dédiées à la compréhension de la structuration des biomolécules du vivant, la \textbf{visualisation moléculaire} est l'un des piliers majeurs pour l'acquisition et la transmission de connaissances, et ce depuis le début de la biologie structurale dans les années 50. Comme pour tout phénomène structurel à observer, il est naturel que le processus d'analyse passe par l'observation de celui-ci. Aucun indicateur quantitatif ne décrivant une particularité géométrique mieux que sa visualisation directe. Elle rencontre cependant ses limites avec la complexité croissante des modèles structuraux à observer. La taille des complexes moléculaires obligent à repenser les méthodes d'exploration usuelles en visualisation moléculaire, basée autour du concept de manipulation.


%%%%%%%%%%%%% RV %%%%%%%%%%%%%%%%%%

Ces 10 dernières années ont vu en effet l'essor de la \textbf{Réalité Virtuelle} comme support interactif de choix pour immerger l'utilisateur au coeur de ses expériences, qu'elles soient ludiques ou professionnelles. Elle offre la possibilité de créer un espace de visualisation 3d et à 360 degrés où l'expert peut s'immerger afin d'explorer ses données sans restrictions d'espace. Nous avons donc décrit dans une seconde partie les principales caractéristiques de la RV, ses concepts, ses supports technologiques ainsi que son évolution, passant à la fois par l'évolution de ses techniques d'immersion (affichage, audio), d'interactions (périphériques, retours sensoriels, gestes) mais également de navigation (suivis de mouvement, périphériques fixes). 

Nous avons cependant mis en avant une orientation importante des évolution de la navigation vers l'exploration de scènes virtuelles réalistes dont les repères spatiaux et les rappels du réels sont nombreux. La visualisation et l'exploration de données scientifiques souffrent elles d'un certain retard dans l'adaptation des paradigmes de navigation les concernant. Il est possible d'expliquer ce retard par celui prit par la RV pour être utilisée comme support d'applications scientifiques. Les développements récents ont pourtant réduit ce retard et de nombreux domaines scientifiques utilisent désormais des EV immersifs pour répondre à certaines problématiques. L'apprentissage et la visualisation de données complexes ont pu par exemple profiter de l'apport de la perception de profondeur et la mise en place de techniques d'interactions proches des outils et instruments du monde réel. Usuellement articulée autour de la manipulation des données, nous avons également souligné que la navigation en conditions immersive ne peut se faire de la même façon qu'au sein des stations de travail habituelles. 


%%%%%%%%%%%% Navigation guidée %%%%%%%%%%%%%

Fort du constat de l'absence de paradigmes spécifiques pour l'exploration de structures moléculaire dans des dispositifs immersifs, nous avons proposé dans une troisième partie plusieurs paradigmes de navigation basés sur le contenu moléculaire et les tâches que les experts scientifiques peuvent être amené à effectuer dans des EV immersifs. Ces paradigmes répondent aux différentes échelles de granularité que présente la visualisation de données scientifiques complexes, à savoir une première approche d'exploration globale avant un focus sur des phénomènes locaux et détaillés. Notre développement a pris en compte au maximum la problématique du malaise du virtuel que nous avons mis en avant comme très présent lors de l'exploration de données abstraites comme peuvent l'être les complexes moléculaires. Nous avons tiré parti de la particularité géométrique commune des complexes moléculaire de grande taille, leur agencement symétrique.

Grâce à l'utilisation d'axes et centres de symétries, nos paradigmes de navigation contraignent l'utilisateur autour de chemins de navigation optimaux pour l'exploration de structures moléculaires de grande taille. Nous avons considérés des chemins de navigation externes et internes aux complexes et mis en place des solutions de génération automatique de chemins dans le cadre de tâches expertes spécifiques comme l'accès à des régions d'intérêt enfouies dans le complexe ou la comparaison de phénomènes biologiques répétés sur les différentes sous-unités consituant le complexe moléculaire.

%%%%%%%%%%% Appli mobile %%%%%%%%%%%%%%%%%


Motivés par les apports de la RV dans les technologies de visualisation avancées et l'approche d'immersion cognitive, notre intérêt s'est ensuite porté sur les moyens de communication dans le monde scientifique. Ces derniers ont connu une évolution rapide ces dernières années profitant de l'évolution des moyens de communication standards. Ces derniers sont maintenant essentiellement basés autour du réseau internet et du support informatique. Au-delà de la simple lecture d'articles scientifiques sur un écran fixe, on tend à déporter cette lecture sur des périphériques mobiles type smartphones ou tablettes, ces derniers possédant la résolution nécessaire pour une expérience de lecture satisfaisante. Les moyens d'interactions avec les informations scientifiques ont également changé, profitant de l'outil informatique pour ajouter des couches de données aux simples textes et images.

Notre approche s'est inscrite dans cette évolution et a pour but de fournir une méthode simple et rapide pour explorer des données de structures 3d de molécules à partir de fichiers légers et facilement échangeables ou téléchargeables. Basée sur un moteur de jeu largement utilisé au sein de la communauté des développeurs d'application mobiles, Unity3D, nous avons créé une application permettant de générer un objet 3d correspondant à une scène moléculaire grâce à l'utilisation de seulement 2 fichiers images, une image de texture et une carte de profondeur. Grâce à l'utilisation du gyroscope présent dans la grande majorité des périphériques mobiles, nous pouvons permettre à l'utilisateur de visualiser sous plusieurs angles une structure moléculaire avec la perception de profondeur qu'offre la 3d.
Dans la continuité de notre approche, nous avons également mis au point la possibilité d'explorer un modèle 3d complet de molécule en situant l'utilisateur au milieu de sa structure et permettant à l'utilisateur d'utiliser son smartphone comme une fenêtre sur le monde virtuel.

%%%%%%%%%%% Visu Ana %%%%%%%%%%%%%%%%

La quatrième partie du manuscrit réintroduit la volonté de mettre l'utilisateur au coeur du contrôle de sa simulation afin de lui permettre de filtrer les données générées et ainsi réduire la quantité de données rapatriée depuis les centres de calcul vers les laboratoires. Nous sommes revenus sur la nécessité d'intégrer un tel schéma de travail afin de faire évoluer le schéma actuel, passant aujourd'hui par une dissociation des étapes de visualisation et d'analyses, pourtant tout deux étroitement liés et en communication constante.
Le rapprochement de ces deux espaces est l'un des sujets d'étude central du \textit{Visual Analytics}. Ce domaine se base sur la mise en place d'une interactivité avancée au coeur de la visualisation de données scientifiques. Notre réflexion autour des techniques de \textit{Visual Analytics} nous ont amené à considéré comme une solution optimale, la réunion des concepts manipulés au sein des espaces de visualisation et d'analyses.
Nous avons montré que grâce aux outils du web sémantique, il est possible de stocker autour d'une architecture définie l'ensemble des données générées au cours d'un processus d'étude de biologie structurale. Au-delà de la simplicité d'accès à ces données, il est possible de raisonner sur celles-ci et de mettre en place des liens forts entre leurs différentes représentations graphiques. Il est ainsi aisé de créer des méthodes d'interactions simplifiées permettant de passer de la visualisation à l'analyse et inversement dans un temps interactif.

Cette fusion des activités n'a pu se faire que grâce à la mise en place d'une description haut-niveau des concepts mis en jeu dans ces deux activités. La cinquième et dernière partie de notre manuscrit met en avant notre réflexion autour des récentes avancées faites dans le Web Sémantique pour parvenir à créer une ontologie métier spécialisée nous permettant de structurer les différents concepts d'analyses et de visualisation et de les lier autour d'une même problématique, la visualisation analytique structurale.
Nous avons donc détaillé la création d'une plateforme de travail mixant visualisation et analyses au sein d'un même espace de travail immersif. Après avoir introduit les différentes facettes de notre plateforme, nous l'avons évalué au moyen de scénario précis et de méthodes ergonomiques existantes.

Conscients de la nécessité d'optimiser les processus d'interactions afin de coller au plus près des standards imposés par la RV et profitant de la capacité d'immersion qu'elle fournit, nous avons cherché à mettre en place un cadre logiciel générique permettant de lier visualisation et analyses au sein d'un même espace de travail.

%%%%%%%%% Sémantique %%%%%%%%%%%%%%



L'intégration complète de ces outils au sein d'un dispositif immersif est le dernier pas que nous devons encore franchir pour proposer un outil complet opérationnel aux experts du domaine et constituerait la seule nuance apportée à l'ensemble de nos contributions.




\section*{Perspectives}

\subsection{Immersion cognitive via périphériques mobiles}

\subsubsection*{Intégration au sein d'un environnement d'utilisation dédié}

L'une des forces de notre application DepthMol3D est sa simplicité d'utilisation couplée à son absence de contraintes matérielles et logicielles. Dans le monde connecté d'aujourd'hui, il est important de permettre un accès aux données à communiquer à tout instant. L'utilisation de technologies de QR-codes est une piste crédible pour permettre d'accélérer le processus d'acquisition des images nécessaires à l'application. Grâce à des méthodes de stockage dans le \textit{cloud}, il est possible d'imaginer la possibilité de rendre disponible à tout instant un ensemble de cartes de profondeur et leurs textures associées pouvant être téléchargées et chargées par l'application suite à la simple détection d'un QR-code (voir Figure \ref{Fig:qr_code}).

\begin{figure}[h]
  \centering
  {\includegraphics[width=.75\linewidth]{./figures/conclusion/qr_code.jpg}}
    \caption{{\it Exemple de QR-code redirigeant vers un site web.}}
  \label{Fig:qr_code}
  \hspace{0.2cm}
\end{figure}

\subsection*{Amélioration des repères spatiaux et élargissement des tâches expertes supportées}

Le retour d'expérience des utilisateurs d'un logiciel est important mais néanmoins souvent oublié quand on arrive aux étapes de post-développement. Avec l'évolution rapide des techniques, plateformes et approches de visualisation, il est important de revoir les nouveaux besoins des experts et les nouvelles possibilités offertes par les nouvelles technologies. 

Il nous semble également important d'utiliser l'expertise de l'utilisateur afin de créer des paradigmes de navigation répondant à leurs besoins, ceux-ci évoluant avec la nature du complexe moléculaire étudié. Ce plus haut niveau d'interaction ne doit pourtant pas dépasser le niveau d'implication cognitive que doit comporter la navigation, cette dernière restant bien souvent un moyen d'accéder à d'autres tâches plus complexes et nécessitant davantage de réflexion de la part de l'utilisateur. Et même si les études ont montré qu'une augmentation de la charge cognitive est une bonne manière de réduire le \textit{cybersickness} (voir section \ref{cybersickness}) elle peut également réduire l'efficacité des tâches annexes au processus de navigation.



\subsection*{Visualisation analytique et sémantique}

\subsubsection{Multimodalité pour l'interaction vocale}

\subsubsection{Automatisation de la génération de graphes}

\section*{Remerciements spécifiques}

“This work was conducted using the Protégé resource, which is supported by grant GM10331601 from the National Institute of General Medical Sciences of the United States National Institutes of Health.” 